\begin{definition}[Joining labels of two actual open regions]
    \label{def:peps_open_region_joining_labels}
    \lean{TNLean.PEPS.RegionJoiningEdge,
      TNLean.PEPS.regionJoiningEdge_ofAdj,
      TNLean.PEPS.regionGluingLabels,
      TNLean.PEPS.regionGluingBoundary}
    \leanok
    Let $L,S$ be disjoint vertex regions of a finite simple graph.
    Their joining bonds are the bonds internal to $L\cup S$ but internal
    to neither part. For joining labels $\mu$ and prescribed outer boundary
    labels $\theta$, let $\beta_L(\mu,\theta)$ and $\beta_S(\mu,\theta)$
    read the original crossing labels of the two parts: they equal $\mu$
    on joining bonds and $\theta$ on the exterior boundary.
    These are the actual labels in the two-column blocking of
    \cite{Schuch2010PEPS}, lines 2489--2507.
\end{definition}

\begin{theorem}[Joining two open contractions]
    \label{thm:peps_graph_open_region_gluing}
    \lean{TNLean.PEPS.regionJoiningEdge_boundary,
      TNLean.PEPS.graphOpenRegionNetwork_union}
    \leanok
    \uses{def:peps_open_region_joining_labels}
    Write $\Psi_R(\theta)$ for the original open tensor contraction on $R$.
    With vertex-dependent physical alphabets and regular group bond labels,
    \[
      \Psi_{L\cup S}(\theta)
        =\sum_\mu \Psi_L(\beta_L(\mu,\theta))
                    \otimes\Psi_S(\beta_S(\mu,\theta)).
    \]
    Every exterior label remains fixed, and each joining label is summed
    exactly once. No factorization identity is assumed.
\end{theorem}
\begin{proof}\leanok
    Split the internal labels of the union into the internal labels of
    $L$, the internal labels of $S$, and the joining labels. The vertex
    product splits over the disjoint union. Summing each internal family
    gives its original open-region contraction.
\end{proof}

\begin{definition}[Permuting accessible boundary registers]
    \label{def:peps_boundary_register_permutation}
    \lean{TNLean.PEPS.regularBoundaryRegisterPermutation}
    \leanok
    For a rooted spanning tree of a connected region, permute its accessible
    boundary registers by $y_b\mapsto y_{\tau b}$ and retain all internal
    reference, normalized vertex and cycle registers. This is the boundary
    exchange in the two independent column operations of
    \cite{Schuch2010PEPS}, lines 2489--2507.
\end{definition}

\begin{theorem}[Original-spin implementation of a boundary exchange]
    \label{thm:peps_regular_boundary_register_transport}
    \lean{TNLean.PEPS.regularBoundaryRegisterPermutation_coordinateTranslation,
      TNLean.PEPS.exists_unitary_regularBoundaryRegisterTransport}
    \leanok
    \uses{def:peps_boundary_register_permutation,
      thm:peps_regular_coordinate_permutation_support,
      thm:peps_regular_physical_unitary_transport}
    For an untwisted connected region with locally regular $G$-isometric
    original sites, one original-spin unitary $W_\tau$ satisfies
    \[
      W_\tau\Psi_R(\theta)=\Psi_R(\theta\circ\tau)
      \qquad\text{for every actual boundary configuration }\theta.
    \]
    The spanning tree is chosen internally, and $W_\tau$ is chosen before
    $\theta$. Its preservation of the full Gram form is derived from the
    original local isometries.
\end{theorem}
\begin{proof}\leanok
    Boundary permutations commute with independent vertex translations
    in accessible coordinates. They therefore commute with the product
    averaging projector. The derived local Gram identity lifts each
    permutation to the original spins; the actual open columns have the
    corresponding permuted boundary labels.
\end{proof}

\begin{theorem}[Literal diagonal weights on a joining bond]
    \label{thm:peps_regular_boundary_charge_contraction}
    \lean{TNLean.PEPS.graphOpenRegionNetwork_regularEdgeCharacterSite_boundary,
      TNLean.PEPS.graphOpenRegionNetwork_regularEdgeCharacterSite_union}
    \leanok
    \uses{thm:peps_graph_open_region_gluing}
    A diagonal insertion $\operatorname{diag}\chi(p\cdot)$ on a joining
    bond $e$ contributes exactly one factor $\chi(p\mu_e)$ to the joined
    contraction. On its tail part this is the prescribed boundary weight;
    on the other part no site factor changes. Thus
    \[
      \Psi_{L\cup S}^{e,\chi,p}(\theta)
        =\sum_\mu\chi(p\mu_e)
          \Psi_L(\beta_L(\mu,\theta))\otimes
          \Psi_S(\beta_S(\mu,\theta)).
    \]
    This is the literal charge insertion of \cite{Schuch2010PEPS},
    Section~6.6, source lines 2432--2453.
\end{theorem}
\begin{proof}\leanok
    The ordered tail occurs in exactly one part. Its scalar character
    factor is read from the common joining label, irrespective of endpoint
    sorting; the other site factors remain the original factors.
\end{proof}

\begin{definition}[The crossing bonds of either part]
    \label{def:peps_joining_boundary_edge}
    \lean{TNLean.PEPS.regionJoiningBoundaryEdge}
    \leanok
    Every actual joining bond is a crossing bond of each of the two
    disjoint regions, with the same underlying native edge.
\end{definition}

\begin{theorem}[Charge motion by independent original-spin operations]
    \label{thm:peps_regular_two_part_charge_motion}
    \lean{TNLean.PEPS.exists_unitary_regularTwoPartChargeMotion}
    \leanok
    \uses{def:peps_joining_boundary_edge,
      thm:peps_regular_boundary_register_transport,
      thm:peps_regular_boundary_charge_contraction}
    Let $L,S$ be disjoint connected regions with locally regular
    $G$-isometric original sites, and let $e,f$ be joining bonds.
    There are independent unitaries $W_L,W_S$, chosen before every
    coefficient function $\chi$, parameter $p$ and exterior configuration
    $\theta$, such that
    \[
      (W_L\otimes W_S)\Psi_{L\cup S}^{e,\chi,p}(\theta)
        =\Psi_{L\cup S}^{f,\chi,p}(\theta).
    \]
    The physical tensor product is identified with the original site
    coordinates by the disjoint-union coordinate equivalence.
\end{theorem}
\begin{proof}\leanok
    Exchange the two joining boundary registers in each part and fix all
    its other crossing registers. The two independent Gram-preserving
    exchanges replace $\mu$ by its common transposition. Relabel this
    finite sum: the single character weight then moves from $e$ to $f$,
    with $\chi,p$ and every external label retained.
\end{proof}

\begin{definition}[The actual two-column torus square]
    \label{def:peps_torus_two_column_charge_geometry}
    \lean{TNLean.PEPS.torusChargeMotionColumn,
      TNLean.PEPS.torusColumnChargeMotionBond,
      TNLean.PEPS.torusColumnChargeMotionBond_val}
    \leanok
    At $v=(x,y)$ the two columns are
    $L=\{(x,y),(x,y+1)\}$ and
    $S=\{(x+1,y),(x+1,y+1)\}$.
    The upper and lower horizontal native bonds join them.
    All coordinates are torus coordinates, including at either seam.
\end{definition}

\begin{theorem}[Derived native column geometry]
    \label{thm:peps_torus_two_column_charge_geometry}
    \lean{TNLean.PEPS.torusChargeMotionColumn_card,
      TNLean.PEPS.torusChargeMotionColumn_disjoint,
      TNLean.PEPS.torusChargeMotionColumn_union,
      TNLean.PEPS.torusChargeMotionColumn_connected}
    \leanok
    \uses{def:peps_torus_two_column_charge_geometry,
      def:peps_torus_plaquette_flux_geometry}
    On a finite torus with both periods at least three, each column contains
    exactly two sites and its actual induced graph is connected. The
    columns are disjoint and their union is the original plaquette region.
\end{theorem}
\begin{proof}\leanok
    The native vertical edge supplies the two-site connecting walk.
    Distinct adjacent horizontal and vertical coordinates give the
    cardinalities and disjointness. The actual plaquette walk has precisely
    these four sites in its support.
\end{proof}

\begin{theorem}[The paper's two-column charge-motion operation on a finite torus]
    \label{thm:peps_torus_two_column_charge_motion}
    \lean{TNLean.PEPS.IsGIsometric.exists_unitary_torusTwoColumnChargeMotion}
    \leanok
    \uses{thm:peps_regular_two_part_charge_motion,
      thm:peps_torus_two_column_charge_geometry}
    Suppose the original four-leg tensor is $G$-isometric for the regular
    action. On every translated square of a finite torus with periods at
    least three, independent two-spin unitaries on its left and right
    columns move the literal diagonal charge from the upper horizontal
    bond to the lower bond. They are chosen before its character, internal
    parameter and every actual outer boundary label, and retain these data.

    This proves the displayed two-column operation of
    \cite[Section~6.6]{Schuch2010PEPS}, source lines 2489--2507, in the stated finite-torus geometry, including both
    seams. Unrestricted lattice geometry remains separate in
    \cite{gap:rmp_peps_quantum_double_g_isometry}.
\end{theorem}
\begin{proof}\leanok
    Use the derived native column geometry and genuine joining bonds in
    the preceding independent-operation theorem. Numbering the group
    labels identifies its coefficient formula with the ordinary open
    contraction of the original charge-inserted site tensors. Each
    column Gram form and the contraction identity are derived internally.
\end{proof}
